% Generated draft: requires mathematical and source review.
\documentclass[11pt]{article}
\usepackage[margin=25mm]{geometry}
\usepackage{fontspec}
\setmainfont{Latin Modern Roman}
\usepackage{amsmath,amssymb,mathtools,unicode-math}
\setmathfont{Latin Modern Math}
\usepackage{enumitem,needspace,xurl,hyperref}
\setlength{\parindent}{0pt}
\setlength{\parskip}{5pt}
\setlength{\emergencystretch}{3em}
\newcommand{\N}{\mathbb{N}}
\newcommand{\Z}{\mathbb{Z}}
\newcommand{\Q}{\mathbb{Q}}
\newcommand{\R}{\mathbb{R}}
\newcommand{\C}{\mathbb{C}}
\newcommand{\F}{\mathbb{F}}
\newcommand{\A}{\mathbb{A}}
\newcommand{\PP}{\mathbb{P}}
\newcommand{\E}{\mathbb{E}}
\newcommand{\eps}{\varepsilon}
\newcommand{\dd}{\,\mathrm{d}}
\begin{document}
\section*{UnsolvedMath ID: 20001305}
\textbf{Source problem code:} catalogue entry
\subsection{Definitions and mathematical statement}
Let $\mathcal K^n$ be the nonempty compact convex subsets of $\R^n$, including singletons, with Hausdorff topology, and let $h_K(u)=\max_{x\in K}\langle x,u\rangle$. Rolf Schneider's question asks to classify continuous $T:\mathcal K^n\to\mathcal K^n$ satisfying
\[
h_{T(K+L)}=h_{T(K)}+h_{T(L)},\qquad T(\theta K)=\theta T(K)\quad(\theta\in SO(n)).
\]
The first identity is Minkowski additivity. Translation invariance is not assumed. Rotations are taken to mean proper rotations, an essential convention in dimension two.

For $n\ge2$, the attached aid gives the unique decomposition
\[
T(K)=\Phi(K)+\{A s(K)\},
\]
where $s(K)=n\int_{S^{n-1}}u h_K(u)\,\dd\sigma(u)$ is the Steiner point for probability-normalized spherical measure, $A$ is linear and commutes with $SO(n)$, and $\Phi$ is a translation-invariant continuous rotation-equivariant Minkowski-additive operator. The map on singletons is $T(\{x\})=\{Ax\}$ and is exactly the translation defect. For $n\ge3$, $A=aI$; for $n=2$, $A=aI+bJ$ with $J$ rotation by $\pi/2$. Requiring reflection equivariance would remove the $bJ$ term, so it cannot be removed under the stated proper-rotation hypothesis alone.

For the normalized part the aid asserts a unique centered signed generating-measure representation
\[
h_{\Phi(K)}=h_K*\mu,
\]
using the appropriate rotation convolution and stabilizer-zonal convention. In dimension two the stabilizer for the $SO(2)$ action is trivial, so no unmentioned reflection symmetry is imposed. The nontrivial admissibility requirement is that the convolution be a support function for every $K$. The summary does not give the precise representation theorem's regularity hypotheses or an intrinsic description of the higher-dimensional generator cone.

Its harmonic audit states that the first multiplier vanishes, $m_1=0$, the constant multiplier satisfies $m_0\ge0$, and $m_0=0$ forces $\Phi=0$. With probability-normalized convolution, $m_0=\mu(S^{n-1})$, and centering removes the first moment. The aid reports full classifications in the planar and weakly monotone settings, while leaving the intrinsic general higher-dimensional cone open, but does not reproduce those classification criteria. These are partial structural reductions of the original classification question, not a complete explicit cone description. The $n=1$ case is outside the supplied decomposition discussion.
\subsection{Short English statement}
Classify continuous rotation-equivariant Minkowski-additive operators without silently adding translation invariance. The supplied aid isolates the singleton translation defect, the planar rotation exception and the normalized convolution problem, whose full higher-dimensional generator cone is not specified.
\subsection{Sources}
\begin{itemize}
\item Local UnsolvedMath registry, permanent ID $20001305$, source code catalogue entry; canonical \texttt{aim-convex-geometry-notes.json}, \texttt{notes[36]}.
\item American Institute of Mathematics, \emph{Fourier Analytic Methods in Convexity}, problem 9 attributed to Rolf Schneider: \url{https://aimath.org/WWN/fourierconvex/fourierconvex.pdf}. A terminal page-number spill is omitted. The supplied verification reports the official wording; no PDF inspection is claimed here. No original literature references are supplied.
\item Attached research attempt 1: partial result, run and registry status partially solved, rigor proved, novelty confidence low, classification PARTIAL-PROGRESS. Its decomposition, representation and multiplier claims are machine generated and not peer reviewed; independently verify before citation. No exact theorem bibliography is provided.
\end{itemize}
\end{document}
